\chapter{Experiments I: The Anomaly-Routing Benchmark}\label{ch:exp1}

This chapter is \emph{breadth} evidence, not the headline. Its job is to show that the
knowability trichotomy of Chapters~\ref{ch:theory1}--\ref{ch:theory2} --- helpful,
harmful, unknowable, with a valid certificate that commits or abstains accordingly ---
is not an artifact of one architecture or one kind of shift, but holds across a wide
span of modalities and detectors where ground-truth benefit is exactly computable. The
deep test-time-adaptation safety certificate of Chapter~\ref{ch:exp2} (CIFAR-10-C,
ImageNet-C, and the natural-shift battery) is the headline application; this chapter
lays the foundation underneath it. We therefore read the results that follow as
\emph{coverage}: a low-cost, high-diversity testbed on which the trichotomy can be
exercised end to end and falsified if it fails, before the deep-network experiments
stress it where it matters most.

Concretely, this chapter evaluates the \KGA\ certificate on a 123-task
anomaly-detection routing benchmark, purpose-built to interrogate the
adapt/freeze/abstain trichotomy of Chapter~\ref{ch:kga} in a domain where ground-truth
benefit is computable, where controlled corruption can inject known regime changes, and
where six classical detectors supply the label-free evidence vector $Z$ without any
neural retraining.  Sections~\ref{sec:exp1-archive}--\ref{sec:exp1-threats} proceed from
data provenance through seven experimental protocols to a cross-experiment synthesis and
a threats-to-validity audit.  Every number is drawn directly from the named JSON result
files in \texttt{experiments/kbound/results/}; discrepancies with \texttt{kbound.tex}
are flagged explicitly in Section~\ref{sec:exp1-threats}.

% ============================================================
\section{The 123-Task Score Archive}
\label{sec:exp1-archive}
% ============================================================

\subsection{Provenance and Construction}

The anomaly-routing benchmark draws tasks from ADBench~\cite{han2022adbench} and the
PyOD library~\cite{zhao2019pyod}, spanning tabular, image-OOD, text, cyber-security, and
fraud-detection datasets.  For each of the 123 tasks we precompute anomaly scores under
six classical detectors: Isolation Forest~\cite{liu2008isolation}, Local Outlier
Factor~\cite{breunig2000lof}, One-class SVM~\cite{scholkopf2001estimating}, and
ECOD~\cite{li2022ecod} form the core; rank-normalised stacking and a reliability-fusion
variant complete the six.  All detector scores are precomputed and cached; no neural
network is retrained at test time.

Each task is split into a \emph{labeled validation set} (used by \KGA\ to estimate $\Delta$
and calibrate the conformal band $\varepsilon$) and a \emph{held-out test set} (used only for
final evaluation).  The canonical frozen candidate $f_0$ is the detector achieving the
best validation AUROC; the adaptation candidate $f_a$ is the rank-normalised logistic
stack in the clean suite, swapped for controlled alternatives in the corrupted regimes of
later experiments.

The labeled validation is the only supervised signal visible to the certificate.  Test
labels are withheld from $Z$ computation: the evidence vector
$Z = (\textit{val\_max}, \textit{val\_gap}, \textit{val\_mean}, \textit{ks\_mean},
\textit{ks\_max}, \textit{disagree}, \textit{val\_std}, \textit{anomaly\_rate},
\textit{log\_ntest})$ is derived entirely from validation performance and distributional
statistics of the unlabeled test batch~\cite{kbound2026repo}.

\subsection{Evaluation Methodology}

For each task $t$ we measure the \emph{true benefit} $B_t = \mathrm{AUROC}(f_a, \mathcal{D}_t)
- \mathrm{AUROC}(f_0, \mathcal{D}_t)$ on the test set.  The four policies are compared by
mean test AUROC:
\begin{itemize}
  \item \textbf{Always-adapt}: apply $f_a$ on all 123 tasks.
  \item \textbf{Always-freeze}: apply $f_0$ on all 123 tasks.
  \item \textbf{\KGA\ (trichotomy)}: apply $f_a$ if $\widehat\Delta - \varepsilon > 0$
    (adapt), apply $f_0$ if $\widehat\Delta + \varepsilon < 0$ (freeze), abstain otherwise.
  \item \textbf{Oracle}: apply whichever of $f_a$ or $f_0$ achieves higher test AUROC.
\end{itemize}
Regret-to-oracle is $r = \mathrm{AUROC}(\text{oracle}) - \mathrm{AUROC}(\text{policy})$;
lower is better.  Safety is reported as adapt precision
$\Pr(B > 0 \mid \text{ADAPT})$ and false-adapt rate $\Pr(B < 0 \mid \text{ADAPT})$.

The conformal level is $\alpha = 0.1$ throughout the clean and harmful suites.  All scripts
are archived in \texttt{kbound\_paper/scripts/} and the repository~\cite{kbound2026repo}.

\subsection{Why Anomaly Routing Is a Clean Testbed for the Trichotomy}

Anomaly detection is structurally well suited for studying the adapt/freeze/abstain
trichotomy for three reasons.  First, benefit $B$ is computable from labeled validation
data, so the ground truth for a decision is unambiguous.  Second, the domain naturally
exhibits the full phase diagram (Chapter~\ref{ch:theory1}): some tasks are
knowably-beneficial, some knowably-harmful, and many are genuinely ambiguous---yielding
a nontrivial trichotomy under any reasonable certificate.  Third, the six-detector score
archive provides a rich, diverse evidence vector $Z$ that is entirely label-free at test
time, directly matching the \KBound\ problem statement (Definition in
Chapter~\ref{ch:formulation}).

The benchmark deliberately does \emph{not} require a powerful learned feature extractor:
the same evidence vector is used across all 123 tasks, and no task-specific tuning is
performed.  This makes the benchmark a conservative test of the certificate; any future
enrichment of $Z$ can only improve results.

% ============================================================
\section{The Knowability Trichotomy Suite}
\label{sec:exp1-trichotomy}
% ============================================================

\subsection{Setup and Protocol}

The candidate adaptation is the rank-normalised logistic stack, a mild ensemble improvement
over the best-validation single detector.  On the clean 123-task suite this candidate
\emph{helps on the majority of tasks}, making it a helpful-dominated regime.  The conformal
band is calibrated at $\alpha = 0.1$; the resulting $\varepsilon$ from
\texttt{knowability\_results.json} is $0.0492$.

\subsection{Decision Distribution}

Table~\ref{tab:exp1-trichotomy} summarises the \KGA\ decisions and safety metrics.

\begin{table}[ht]
\centering
\caption{\KGA\ decisions and safety on the clean 123-task suite
  (source: \texttt{knowability\_results.json}, field
  \texttt{metrics}).}
\label{tab:exp1-trichotomy}
\begin{tabular}{lcc}
\toprule
Metric & Value & Source field \\
\midrule
Tasks total & 123 & \texttt{n\_tasks} \\
Conformal $\varepsilon$ & 0.0492 & \texttt{eps\_conformal} \\
Coverage (acted = adapt or freeze) & 0.171 & \texttt{coverage} \\
Decisions: ADAPT & 20 & \texttt{n\_adapt} \\
Decisions: FREEZE & 1 & \texttt{n\_freeze} \\
Decisions: ABSTAIN & 102 & \texttt{n\_abstain} \\
Adapt precision $\Pr(B>0\mid\text{ADAPT})$ & 0.90 & \texttt{adapt\_precision\_(B>0|ADAPT)} \\
False-adapt rate $\Pr(B<0\mid\text{ADAPT})$ & 0.10 & \texttt{false\_adapt\_rate\_(B<0|ADAPT)} \\
Mean $|\Delta|$ on abstained tasks & 0.0215 & \texttt{abstain\_mean\_|B|} \\
Mean $|\Delta|$ on acted tasks & 0.1324 & \texttt{nonabstain\_mean\_|B|} \\
\midrule
Mean AUROC: \KGA\ (trichotomy) & 0.7661 & \texttt{policy.mean\_auc\_trichotomy} \\
Mean AUROC: always-adapt & 0.7771 & \texttt{policy.mean\_auc\_always\_adapt} \\
Mean AUROC: always-freeze & 0.7481 & \texttt{policy.mean\_auc\_always\_freeze} \\
Mean AUROC: oracle & 0.7828 & \texttt{policy.mean\_auc\_oracle} \\
Regret (\KGA\ vs oracle) & 0.0167 & \texttt{policy.regret\_trichotomy\_vs\_oracle} \\
Regret (always-adapt vs oracle) & 0.0057 & \texttt{policy.regret\_always\_adapt\_vs\_oracle} \\
\bottomrule
\end{tabular}
\end{table}

\subsection{Adapt Precision and Safety}

The certificate achieves adapt precision $0.90$: of the 20 tasks on which \KGA\ commits to
adaptation, 18 have $B > 0$.  The false-adapt rate is $0.10$.  This is the \emph{clean}
helpful-dominated regime, so caution is warranted in interpreting the coverage figure: at
$17.1\%$ coverage, only 21 of 123 tasks are acted upon, while 102 are abstained.

\subsection{Abstain-Where-Benefit-Is-Small: Safety Concentration}

The key safety property of the trichotomy is not just that adapt decisions are accurate but
that \emph{abstentions concentrate on tasks where the true benefit is near zero}.  The JSON
confirms this directly: the mean absolute benefit on abstained tasks is $|\Delta|_\text{abstain}
= 0.0215$, while on acted tasks it is $|\Delta|_\text{act} = 0.1324$---a ratio of
$0.1324 / 0.0215 \approx 6.2$.  Abstaining carries only small expected cost per task
($0.0215$ AUROC gap), whereas every acted task was worth $0.1324$ on average.  This is the
empirical manifestation of the certificate's one-sided guarantee (Theorem~\ref{thm:cert}).

\subsection{Interpretation: Helpful-Dominated Regime}

An honest reading of this suite is that the clean benchmark is helpful-dominated: the
rank-normalised stack mostly helps, so always-adapt (AUROC $0.7771$) outperforms the safe
\KGA\ policy (AUROC $0.7661$) by $0.011$.  The regret gap is $0.0057$ vs $0.0167$---\KGA\
carries roughly three times more regret than always-adapt in this regime.  This is not a
failure: the certificate is designed to be safe, not to maximise AUROC when adaptation is
universally beneficial.  The harmful-regime and mixed-regime experiments
(Sections~\ref{sec:exp1-harmful} and~\ref{sec:exp1-mixed}) demonstrate where the
trichotomy earns its advantage.

% ============================================================
\section{The Harmful-Fusion Regime}
\label{sec:exp1-harmful}
% ============================================================

\subsection{Setup and Protocol}

The harmful-fusion experiment evaluates \KGA\ against a candidate that is \emph{adversarially
bad}: the legacy reliability-fusion candidate (\texttt{legacy\_fuse}; artifact key
\texttt{elara\_fuse}) that hurts on $79.7\%$
of the 123 tasks.  The true harmful base rate, from field
\texttt{base\_rate\_harmful\_B\textless{}0} in \texttt{kbound\_harmful\_results.json}, is
$0.7967$.  The frozen candidate $f_0$ is again the auto-selected best-validation detector
(\texttt{auto\_select}).  If \KGA\ cannot identify and block this harmful fusion, it fails
the most important safety test.

\subsection{Results}

Table~\ref{tab:exp1-harmful} gives the full policy comparison.

\begin{table}[ht]
\centering
\caption{Harmful-fusion regime: AUROC and regret-to-oracle for four policies across 123
  tasks.  Source: \texttt{kbound\_harmful\_results.json}.}
\label{tab:exp1-harmful}
\begin{tabular}{lcccc}
\toprule
Policy & Mean AUROC & Regret vs oracle & Source field \\
\midrule
Always-freeze ($f_0$) & 0.7481 & 0.0018 & \texttt{mean\_auc.always\_freeze(auto\_select)} \\
Always-adapt (legacy-fuse) & 0.7284 & 0.0214 & \texttt{mean\_auc.always\_adapt(elara\_fuse)} \\
\KGA\ (trichotomy) & 0.7479 & 0.0019 & \texttt{mean\_auc.K-Bound\_trichotomy} \\
Oracle & 0.7499 & 0.0000 & \texttt{mean\_auc.oracle} \\
\bottomrule
\end{tabular}
\end{table}

The finding is clean: \KGA\ cuts regret-to-oracle from $0.0214$ (always-adapt) to
$0.0019$, a factor of $0.0214 / 0.0019 \approx 11.3$, and its AUROC ($0.7479$) is
indistinguishable from both always-freeze ($0.7481$) and the oracle ($0.7499$).  The
certificate's job here is to \emph{match} the safe baseline by refusing a dangerous
candidate, not to beat both policies --- always-freeze is already near-oracle when
adaptation almost always hurts.

The decision distribution reinforces the picture: of 123 tasks, \KGA\ issues
\textbf{ADAPT} on 1, \textbf{FREEZE} on 13, and \textbf{ABSTAIN} on 109
(source: \texttt{counts}).  The single adapt decision has $B > 0$ (false-adapt rate
$0.0$ on freezed-or-abstained; field \texttt{freeze\_correct\_B\textless{}=0} $= 1.0$),
and all 13 freeze decisions are on tasks where $B \le 0$ (field
\texttt{freeze\_correct\_B\textless{}=0} $= 1.0$).  The certificate almost entirely
refuses the harmful fusion candidate, correctly identifying the regime.

\begin{figure}[ht]
\centering
\includegraphics[width=0.70\textwidth]{fig_kbound_harmful.png}
\caption{\KGA\ in the harmful-fusion regime.  The certificate almost entirely blocks the
  legacy-fuse candidate; regret-to-oracle falls from $0.0214$ (always-adapt) to $0.0019$
  (${\approx}11\times$ reduction).  Source: \texttt{kbound\_harmful\_results.json}.}
\label{fig:exp1-harmful}
\end{figure}

\subsection{Interpretation}

This experiment directly probes the hard end of the phase diagram: the domain sits deep
in the ``knowably harmful'' region of Chapter~\ref{ch:theory1}'s phase diagram, and \KGA\
correctly refuses to commit.  The $11\times$ regret reduction relative to always-adapt
demonstrates that the certificate's value is largest precisely when the adaptation candidate
is dangerous.  Conversely, always-freeze is already near-oracle in this regime, so the
certificate's main job is to \emph{match} the safe baseline rather than to improve upon it.

\subsection{Threats Specific to This Sub-experiment}

The legacy-fuse candidate is constructed to be harmful; a practitioner would not ordinarily
deploy a candidate with a known $79.7\%$ harmful rate.  This sub-experiment is a
\emph{controlled stress test}, not a claim about typical real-world adaptation candidates.

% ============================================================
\section{The Mixed Regime}
\label{sec:exp1-mixed}
% ============================================================

\subsection{Construction}

The mixed-regime benchmark is this chapter's heterogeneity probe: a single evaluation that
deliberately places three qualitatively different conditions side by side to ask whether
the trichotomy routes by true regime.  We read its AUROC contrast as a diagnostic rather
than a beats-both --- the conference mixed-stream comparison was withdrawn pending an
out-of-fold re-run, as the Policy AUROC Comparison below makes explicit --- while the
decision-distribution and false-adapt findings are what it contributes to the safety
story.  It is built to place three qualitatively different conditions in the same
evaluation:
\begin{enumerate}
  \item \textbf{Clean} ($n_\text{clean} = 123$ instances): the best-validation detector is
    used as $f_a$; adaptation is mildly \emph{harmful} (mean $B = -0.041$, field
    \texttt{true\_mean\_benefit\_B\_by\_condition.clean}).
  \item \textbf{Detectable failure} ($n_\text{det} = 123$ instances): the best-validation
    detector is corrupted with a distribution shift; adaptation is strongly
    \emph{beneficial} (mean $B = +0.193$, field
    \texttt{true\_mean\_benefit\_B\_by\_condition.detectable}) and the shift is observable
    via KS statistics in $Z$.
  \item \textbf{Covert failure} ($n_\text{cov} = 123$ instances): the best-validation
    detector is permuted so its signal is destroyed but the marginal distribution of $Z$
    is nearly unchanged; adaptation is also \emph{beneficial} (mean $B = +0.181$, field
    \texttt{true\_mean\_benefit\_B\_by\_condition.covert}), yet evidence is confounded
    with clean.
\end{enumerate}
Total instances: $3 \times 123 = 369$ (field \texttt{n\_instances}).  The conformal
$\varepsilon = 0.1315$ (field \texttt{eps}) is larger than in the clean suite because
the mixed dataset has higher within-batch variance.

\subsection{Decision Distribution by True Regime}

Table~\ref{tab:exp1-mixed-decisions} shows \KGA\ decisions broken out by condition.

\begin{table}[ht]
\centering
\caption{\KGA\ decisions by true regime in the controlled mixed benchmark.
  Source: \texttt{mixed\_regime\_results.json},
  field \texttt{decision\_distribution\_by\_condition}.}
\label{tab:exp1-mixed-decisions}
\begin{tabular}{lrrrr}
\toprule
Condition & True mean $\Delta$ & ADAPT & FREEZE & ABSTAIN \\
\midrule
Clean            & $-0.041$ & 8  & 5  & 110 \\
Detectable failure & $+0.193$ & 71 & 3  & 49  \\
Covert failure     & $+0.181$ & 59 & 3  & 61  \\
\midrule
Total & & 138 & 11 & 220 \\
\bottomrule
\end{tabular}
\end{table}

The pattern is informative.  On the clean (mildly harmful) condition, \KGA\ abstains on
110 of 123 instances and adapts on only 8---correctly exercising restraint.  On the
detectable-failure (beneficial) condition, it adapts on 71 and abstains on 49---capturing
most of the large-$B$ instances.  On the covert-failure condition, it adapts on 59 and
abstains on 61---split because the permutation partially destroys the evidence signal that
would distinguish this condition from clean.

\subsection{Policy AUROC Comparison}

Table~\ref{tab:exp1-mixed-auroc} gives the four policy AUROC values.

\begin{table}[ht]
\centering
\caption{Policy mean AUROC in the mixed regime (369 instances).
  Source: \texttt{mixed\_regime\_results.json},
  field \texttt{mean\_auc\_policies}.}
\label{tab:exp1-mixed-auroc}
\begin{tabular}{lccc}
\toprule
Policy & Mean AUROC & Regret vs oracle & Source field \\
\midrule
Always-freeze & 0.5862 & 0.1327 & \texttt{always\_freeze} \\
\KGA           & 0.6899 & 0.0291 & \texttt{K\_Bound} \\
Always-adapt  & 0.6971 & 0.0218 & \texttt{always\_adapt} \\
Oracle        & 0.7189 & 0.0000 & \texttt{oracle} \\
\bottomrule
\end{tabular}
\end{table}

In this single in-fold run \KGA\ records a mean AUROC of $0.6899$, above always-freeze
($0.5862$) and within $0.0072$ of always-adapt.  We treat this contrast as a
\emph{diagnostic} observation rather than a claimed beats-both: the conference version of
this work withdrew the mixed-stream result pending a corrected out-of-fold re-run, and we
preserve that caution here.  The gap to always-freeze is consistent with the trichotomy
routing away from the mildly harmful clean condition, but the figure is computed in-fold
on a constructed stream, so we do not advance it as a clean win over both trivial
policies.  The load-bearing safety reading is the same regardless of the AUROC contrast,
and it does not depend on the withdrawn comparison: adapt precision $= 0.935$ and
false-adapt rate $= 0.065$ (fields \texttt{safety.adapt\_precision\_B>0} and
\texttt{safety.false\_adapt\_rate\_B\textless{}0}), against the budget $\alpha = 0.1$.

\begin{figure}[ht]
\centering
\includegraphics[width=0.49\textwidth]{fig_mixed_policies.png}\hfill
\includegraphics[width=0.49\textwidth]{fig_mixed_decisions.png}
\caption{Mixed regime (in-fold, single run; the conference mixed-stream comparison was
  withdrawn pending an out-of-fold re-run, so the AUROC contrast is read as a diagnostic,
  not a beats-both).  \textbf{Left}: policy AUROC comparison---\KGA\ sits above
  always-freeze and within $0.0072$ of always-adapt.  \textbf{Right}: \KGA\ decision map
  by true condition---adapt concentrated on detectable-failure instances, abstain
  dominant on clean.  Source: \texttt{mixed\_regime\_results.json}.}
\label{fig:exp1-mixed}
\end{figure}

\subsection{Non-identifiability on the Clean/Covert Boundary}

An important structural observation from the mixed regime is that the evidence vectors of
clean and covert-failure instances are geometrically close.  The mean $Z$-distance from a
clean instance to its nearest covert-failure neighbour is $1.354$ (field
\texttt{nonidentifiability\_clean\_vs\_covert.mean\_Z\_L2\_distance\_clean\_to\_nearest\_covert}),
while the mean distance to the nearest \emph{detectable}-failure neighbour is $3.100$
(field \texttt{mean\_Z\_L2\_distance\_clean\_to\_nearest\_detectable})---a factor of
$3.100 / 1.354 \approx 2.3$ wider separation.

This is the empirical signature of the partial non-identifiability discussed in
Theorem~\ref{thm:imp}: permutation preserves marginals but only slightly perturbs
correlation structure, so $Z$ is not perfectly identical across worlds yet cannot
reliably distinguish them.  The certificate correctly responds by abstaining more on
the covert-failure condition (61 abstentions) than on the detectable-failure condition
(49 abstentions), even though both have positive true benefit---hedging on the harder
identification problem.

\subsection{Interpretation}

The mixed regime probes whether \KGA\ can \emph{route by true regime} when conditions are
heterogeneous, and the decision distribution by condition (Table~\ref{tab:exp1-mixed-decisions})
is the part of the picture we stand behind: the trichotomy abstains on the mildly harmful
clean condition, concentrates adaptation on the detectable-failure instances, and hedges
on the covert condition.  The accompanying AUROC contrast we treat only as a diagnostic,
for the reason given above --- the mixed-stream comparison was withdrawn pending an
out-of-fold re-run --- and we do not claim it as a beats-both.  What does not rest on that
contrast is the safety result: the false-adapt rate of $0.065$ --- only $8$ of 138 adapt
decisions land on truly harmful clean instances --- is a material improvement over
always-adapt's implicit $100\%$ false-adapt rate on the clean condition, and it is the
property the certificate is built to deliver.

% ============================================================
\section{Statistical Rigor: Eight Seeds and Paired $t$-Tests}
\label{sec:exp1-rigor}
% ============================================================

\subsection{Protocol}

To verify that the mixed-regime results are not an artefact of a single random seed, we
repeat the entire experiment over 8 independent seeds.  Each seed controls the
conformal calibration split and the construction of the evidence vector $Z$.
We use split-conformal calibration at $\alpha = 0.1$ throughout.  All 8 seeds and
their per-seed AUROC values are summarised in Table~\ref{tab:exp1-rigor-agg}, and the
aggregate statistics appear in Table~\ref{tab:exp1-rigor-agg}.

\subsection{Per-Seed Aggregate Results}

\begin{table}[ht]
\centering
\caption{Multi-seed statistics (8 seeds, mixed regime).  Values are mean $\pm$ std.
  Source: \texttt{rigor\_multiseed.json}, field \texttt{mean\_std}.}
\label{tab:exp1-rigor-agg}
\begin{tabular}{lccc}
\toprule
Policy / Metric & Mean & Std & Source field \\
\midrule
\KGA\ AUROC        & 0.6860 & 0.0033 & \texttt{K\_Bound} \\
Always-adapt AUROC & 0.6970 & 0.0001 & \texttt{always\_adapt} \\
Always-freeze AUROC & 0.5842 & 0.0015 & \texttt{always\_freeze} \\
Oracle AUROC        & 0.7180 & 0.0006 & \texttt{oracle} \\
\midrule
Adapt precision    & 0.9358 & 0.0101 & \texttt{adapt\_precision} \\
False-adapt rate   & 0.0642 & 0.0101 & \texttt{false\_adapt} \\
Coverage           & 0.3726 & 0.0175 & \texttt{coverage} \\
\bottomrule
\end{tabular}
\end{table}

The standard deviations are small: \KGA's AUROC varies by only $\pm 0.003$ across seeds,
confirming that the mixed-regime result is stable.  Always-adapt is nearly deterministic
($\text{std} = 0.0001$) because it never abstains.  The adapt precision and false-adapt
rate are complementary ($0.9358 + 0.0642 = 1.000$) as required.

\subsection{Paired $t$-Tests}

We compute paired $t$-tests across the 8 seeds (each seed provides one AUROC observation
per policy).  Results from field \texttt{paired\_ttest\_KBound\_vs\_always\_freeze} and
\texttt{paired\_ttest\_KBound\_vs\_always\_adapt}:

\begin{itemize}
  \item \textbf{\KGA\ vs always-freeze}: $t = 62.24$, $p = 7.26 \times 10^{-11}$.
    Across the eight in-fold seeds \KGA\ scores above always-freeze with low
    between-seed variance.
  \item \textbf{\KGA\ vs always-adapt}: $t = -8.76$, $p = 5.08 \times 10^{-5}$.
    \KGA\ scores below always-adapt; in this mild-harm domain the aggressive policy
    remains ahead.
\end{itemize}

Both signs are reported honestly, and both inherit the caveat of
Section~\ref{sec:exp1-mixed}: these are in-fold seeds of the mixed-stream construction
whose headline comparison the conference version withdrew pending an out-of-fold re-run,
so we read the first test as a stable diagnostic --- the seed-to-seed variance is small
--- rather than as a confirmed beats-freeze claim.  The second test stands on its own
terms: in the mixed regime as constructed, where clean-condition harm is mild
($B = -0.041$) and covert-failure provides genuine benefit, always-adapt is not dominated
by \KGA.  The deep-TTA stress grid of Chapter~\ref{ch:exp2} is the setting that carries
the corruption-grid beats-both claim; this chapter's contribution to the safety story is
the controlled false-adapt rate, which the multi-seed run confirms is stable
($0.0642 \pm 0.0101$).

% ============================================================
\section{Ablations}
\label{sec:exp1-ablations}
% ============================================================

\subsection{Evidence-Component Knockout}

We assess the contribution of each component of the evidence vector $Z$ by performing
leave-one-out ablations: we drop one feature at a time and measure the resulting regret
on the mixed regime.  The baseline (all features) has regret $0.0365$ and false-adapt
rate $0.0650$ (source: \texttt{ablations.json}, field \texttt{baseline}).

Table~\ref{tab:exp1-abl-drop} reports the full ablation results.

\begin{table}[ht]
\centering
\caption{Evidence-drop ablation: regret and false-adapt rate when each $Z$ feature is
  removed.  Source: \texttt{ablations.json}, field \texttt{evidence\_drop}.
  Baseline (all features): regret $0.0365$, false-adapt $0.065$.
  The most important feature is \texttt{disagree}; removing it raises regret from
  $0.0365$ to $0.0940$ ($2.6\times$) and false-adapt from $0.065$ to $0.250$.}
\label{tab:exp1-abl-drop}
\begin{tabular}{lccc}
\toprule
Feature dropped & Regret & False-adapt & Coverage \\
\midrule
\texttt{val\_max}      & 0.0350 & 0.064 & 0.360 \\
\texttt{val\_gap}      & 0.0372 & 0.067 & 0.336 \\
\texttt{val\_mean}     & 0.0426 & 0.071 & 0.320 \\
\texttt{ks\_mean}      & 0.0361 & 0.066 & 0.352 \\
\texttt{ks\_max}       & 0.0373 & 0.066 & 0.352 \\
\textbf{\texttt{disagree}} & \textbf{0.0940} & \textbf{0.250} & 0.141 \\
\texttt{val\_std}      & 0.0365 & 0.074 & 0.344 \\
\texttt{anomaly\_rate} & 0.0366 & 0.065 & 0.350 \\
\texttt{log\_ntest}    & 0.0393 & 0.068 & 0.336 \\
\midrule
Baseline (all) & 0.0365 & 0.065 & --- \\
\bottomrule
\end{tabular}
\end{table}

\begin{figure}[ht]
\centering
\includegraphics[width=0.70\textwidth]{fig_ablation_evidence.png}
\caption{Evidence-component ablation: regret when each feature of $Z$ is removed.
  The \texttt{disagree} feature dominates: removing it raises regret $2.6\times$.
  Source: \texttt{ablations.json}.}
\label{fig:exp1-ablation}
\end{figure}

\paragraph{Disagreement is the most informative feature.}  Removing \texttt{disagree}
raises regret from $0.0365$ to $0.0940$---a $2.6\times$ increase---and the false-adapt
rate from $0.065$ to $0.250$.  This directly supports Theorem~\ref{thm:disagree}
(Chapter~\ref{ch:theory2}): the disagreement signal between $f_0$ and $f_a$ is the
primary evidence of whether adaptation is safe.  All other features cause much smaller
degradations, with \texttt{val\_mean} second at regret $0.0426$.

\subsection{Alpha Sweep: Safety-Coverage Tradeoff}

The conformal level $\alpha$ controls the width of the certificate band $\varepsilon$:
lower $\alpha$ means tighter bands and fewer commits (lower coverage, higher safety);
higher $\alpha$ means wider commits (higher coverage, potentially higher false-adapt).
Table~\ref{tab:exp1-alpha} reports the sweep.

\begin{table}[ht]
\centering
\caption{$\alpha$ sweep: coverage and false-adapt rate at four conformal levels.
  Source: \texttt{ablations.json}, field \texttt{alpha\_sweep}.}
\label{tab:exp1-alpha}
\begin{tabular}{lcccc}
\toprule
$\alpha$ & $\varepsilon$ & Coverage & False-adapt & Source \\
\midrule
0.01 & 0.3240 & 0.114 & 0.000 & \texttt{alpha\_sweep."0.01"} \\
0.05 & 0.1912 & 0.263 & 0.074 & \texttt{alpha\_sweep."0.05"} \\
0.10 & 0.1391 & 0.350 & 0.065 & \texttt{alpha\_sweep."0.1"} \\
0.20 & 0.0910 & 0.488 & 0.072 & \texttt{alpha\_sweep."0.2"} \\
\bottomrule
\end{tabular}
\end{table}

At $\alpha = 0.01$ the certificate achieves $0\%$ false-adapt at $11.4\%$ coverage.
At $\alpha = 0.20$ coverage rises to $48.8\%$ with a false-adapt rate of $7.2\%$.  The
canonical $\alpha = 0.10$ setting ($35\%$ coverage, $6.5\%$ false-adapt) sits near the
elbow of the safety-coverage curve, offering a reasonable operating point.

\subsection{Batch-Size Sweep}

The certificate width $\varepsilon$ depends on the size of the unlabeled test batch:
larger batches yield tighter bands.  Table~\ref{tab:exp1-batch} confirms the theoretical
prediction of Chapter~\ref{ch:theory2}: $\varepsilon$ decreases monotonically from
$0.443$ at batch size $16$ to $0.182$ at batch size $256$, while coverage increases from
$2\%$ to $26\%$.

\begin{table}[ht]
\centering
\caption{Batch-size sweep: $\varepsilon$ and coverage at five batch sizes.
  Source: \texttt{ablations.json}, field \texttt{batchsize\_sweep}.}
\label{tab:exp1-batch}
\begin{tabular}{lccc}
\toprule
Batch size & $\varepsilon$ & Coverage & Source \\
\midrule
16  & 0.4434 & 0.020 & \texttt{batchsize\_sweep."16"} \\
32  & 0.3549 & 0.036 & \texttt{batchsize\_sweep."32"} \\
64  & 0.2777 & 0.199 & \texttt{batchsize\_sweep."64"} \\
128 & 0.2554 & 0.209 & \texttt{batchsize\_sweep."128"} \\
256 & 0.1823 & 0.261 & \texttt{batchsize\_sweep."256"} \\
\bottomrule
\end{tabular}
\end{table}

The coverage gain from $16$ to $256$ samples ($2.0\% \to 26.1\%$) reflects the
$1/\sqrt{n}$ shrinkage of the conformal quantile.  At small batches the certificate is
necessarily conservative; this is a structural limitation acknowledged in
Chapter~\ref{ch:theory2} and the threats section below.

\subsection{Interpretation}

Taken together, the ablations show that (a) \texttt{disagree} is irreplaceable---its
removal degrades safety catastrophically---consistent with Theorem~\ref{thm:disagree};
(b) the remaining features are modestly useful individually but collectively stabilise the
estimate; and (c) both $\alpha$ and batch size trade off smoothly, confirming the
theoretical picture.

% ============================================================
\section{Regression Under Covariate Shift}
\label{sec:exp1-regression}
% ============================================================

\subsection{Setup}

To demonstrate that the \KBound\ certificate extends beyond anomaly detection, we
construct a synthetic regression task.  The data-generating process is $Y = X^\top \beta
+ \epsilon$ with $X \in \R^{10}$, where the source and target distributions satisfy
$P_S(Y \mid X) = P_T(Y \mid X)$ (label shift absent) but have covariate shift up to
$4\sigma$.  The frozen candidate $f_0$ is ordinary least squares (OLS); the adaptation
candidate $f_a$ is importance-weighted least squares (IW), using density-ratio weights
estimated from the label-free test batch.

Importance weighting \emph{hurts} on $79.2\%$ of the 72 test instances (field
\texttt{harmful\_rate\_B\textless{}0} in \texttt{regression\_covariate.json}): high
weight variance inflates the IW estimator's variance enough that it exceeds its bias
correction benefit.  The true harmful rate ($79.2\%$) is comparable to the harmful-fusion
regime ($79.7\%$), making this an equally severe safety test.

\subsection{Results}

\begin{table}[ht]
\centering
\caption{Regression covariate-shift experiment: MSE and \KGA\ decision distribution.
  Source: \texttt{regression\_covariate.json}.
  Lower MSE is better; benefit $B = \text{MSE}(f_0) - \text{MSE}(f_a)$.}
\label{tab:exp1-regression}
\begin{tabular}{lcc}
\toprule
Policy & Mean MSE & Source field \\
\midrule
Always-adapt (IW) & 0.3376 & \texttt{mean\_MSE.always\_adapt(IW)} \\
Always-freeze (OLS) & 0.2513 & \texttt{mean\_MSE.always\_freeze} \\
\KGA & 0.2513 & \texttt{mean\_MSE.K\_Bound} \\
Oracle & 0.2512 & \texttt{mean\_MSE.oracle} \\
\midrule
\multicolumn{3}{l}{\textit{\KGA\ decisions:} ADAPT = 0, FREEZE = 11, ABSTAIN = 61} \\
\bottomrule
\end{tabular}
\end{table}

\begin{figure}[ht]
\centering
\includegraphics[width=0.60\textwidth]{fig_regression.png}
\caption{Regression covariate shift: \KGA\ matches always-freeze and the oracle (MSE
  $0.2513$) by refusing the harmful IW candidate on all 72 instances.
  Always-adapt (IW) reaches MSE $0.3376$.
  Source: \texttt{regression\_covariate.json}.}
\label{fig:exp1-regression}
\end{figure}

\KGA\ issues \textbf{zero adapt decisions}, 11 freeze decisions, and 61 abstain decisions
(source: \texttt{decision\_counts}).  The resulting mean MSE is $0.2513$---matching
always-freeze ($0.2513$) and the oracle ($0.2512$) to within rounding, while always-adapt
(IW) reaches $0.3376$.  Because every instance is either frozen or abstained, the
adapt-precision and false-adapt-rate fields are \texttt{null} (no adapt decisions were
issued; field \texttt{adapt\_precision\_B>0} is \texttt{null}).

\subsection{Connection to Theorem~\ref{thm:pos}}

This experiment instantiates the positive-transfer regime of Theorem~\ref{thm:pos}
(Chapter~\ref{ch:theory2}): the certificate correctly identifies that the density-ratio
variance of the IW weights is too large to trust, and abstains or freezes accordingly.
The result shows that the certificate's safety guarantee transfers \emph{beyond anomaly
detection} to a regression setting without any change to the \KGA\ algorithm, requiring
only that the evidence vector $Z$ includes distributional statistics of the test batch
(here, KS statistics and estimated density-ratio variance).

\subsection{Interpretation}

The regression sub-experiment is deliberately modest in scope (72 instances, synthetic
DGP) but it makes a clean point: the certificate is not anomaly-specific.  The zero
adapt rate is consistent with the high harmful base rate ($79.2\%$), and the oracle match
($0.2512 \approx 0.2513$) confirms that abstaining is the right call when no instance
is clearly adaptable.

% ============================================================
\section{The Clean Non-identifiability Witness}
\label{sec:exp1-witness}
% ============================================================

\subsection{Setup}

The clean witness experiment provides the theory-experiment bridge for
Theorem~\ref{thm:imp} (Chapter~\ref{ch:theory1}).  We construct two worlds:
\begin{itemize}
  \item \textbf{World 0}: $X \sim \mathcal{N}(0,1)$, $f_0 = \ind[X > 0]$,
    $f_a = \ind[X < 0]$, with labels $Y = f_0(X)$.  Adaptation hurts: mean $B = -1.0$
    (field \texttt{world0\_mean\_B}).
  \item \textbf{World 1}: same $X$, $f_0$, $f_a$, but labels $Y = f_a(X)$.  Adaptation
    helps: mean $B = +1.0$ (field \texttt{world1\_mean\_B}).
\end{itemize}
The key property is that \emph{every feature of $Z$ is a function of $X$ only}---it does
not depend on which world is active.  Therefore the law of $Z$ is identical in both
worlds (the marginal of $X$ is the same), while the sign of $B$ is completely opposed.

We draw $M = 300$ instances (field \texttt{M}), interleaving the two worlds, and present
each instance to \KGA.  The conformal band $\varepsilon = 1.3952$ (field \texttt{eps}) is
calibrated from the pooled calibration set.

\subsection{Results}

The KS test on each feature of $Z$ between worlds gives $p$-values all above $0.05$
(field \texttt{all\_Z\_features\_p>0.05} is \texttt{true}):

\begin{itemize}
  \item \texttt{frac\_pos}: KS $p = 0.965$
  \item \texttt{std}: KS $p = 0.942$
  \item \texttt{ks\_vs\_ref}: KS $p = 0.799$
  \item \texttt{mean\_f0}: KS $p = 0.965$
  \item \texttt{skew}: KS $p = 0.165$
\end{itemize}
(Source: \texttt{witness\_clean.json}, field
\texttt{Z\_indistinguishable\_KS\_pvalues}.)

All features are statistically indistinguishable between the two worlds at any conventional
significance level.  \KGA\ responds by \textbf{abstaining on $100\%$} of the 300 instances
(field \texttt{abstain\_rate} $= 1.0$).

The mean $|\widehat\Delta|$ across instances is $0.264$ (field
\texttt{mean\_|\^Bhat|}), but $\varepsilon = 1.395$ dwarfs it: the band covers zero in
every case, so no commit is issued.  A classifier forced to commit by
$\sign\widehat\Delta$ pays regret $0.513$ (field
\texttt{forced\_commit\_regret\_vs\_oracle})---near chance---while the abstaining policy's
regret-to-oracle (defaulting to $f_0$ when abstaining) is $0.467$ (field
\texttt{abstain\_regret\_vs\_oracle(default\_freeze)}).

\subsection{Interpretation: Empirical Realization of Theorem~\ref{thm:imp}}

Theorem~\ref{thm:imp} (the impossibility theorem, Chapter~\ref{ch:theory1}) states
that when the evidence $Z$ has the same law in both the helpful and harmful worlds, no
rule can reliably distinguish them and the optimal action is to abstain.  This experiment
is the exact empirical realisation of that theorem: identical $Z$-law (KS $p > 0.05$ on
all five features), opposite true benefit ($-1.0$ and $+1.0$), and $100\%$ abstention
from \KGA.  The forced-commit regret ($0.513 \approx 0.5$) confirms that committing is
essentially chance-level in this scenario.

This is not a pathological edge case.  It is a deliberately constructed witness to the
fundamental information-theoretic boundary that \KBound\ theory (Theorem~\ref{thm:imp})
identifies.  Real deployments will rarely encounter this exact construction, but the
witness demonstrates that the certificate's abstain behaviour is not merely conservative
risk-aversion---it is information-theoretically necessary.

% ============================================================
\section{Cross-Experiment Regret Synthesis}
\label{sec:exp1-synthesis}
% ============================================================

\subsection{Summary Table}

Table~\ref{tab:exp1-regret-summary} collects the regret-to-oracle numbers from all six
primary sub-experiments of this chapter.

\begin{table}[ht]
\centering
\caption{Cross-experiment regret-to-oracle synthesis.  All figures from the named JSON
  files; the ``vs always-freeze'' column expresses always-freeze's regret for comparison.}
\label{tab:exp1-regret-summary}
\begin{tabular}{llccc}
\toprule
Experiment & Regime & \KGA\ regret & Always-adapt regret & Always-freeze regret \\
\midrule
Trichotomy (clean) & Helpful-dominated & 0.0167 & 0.0057 & 0.0348 \\
Harmful fusion & Harmful-dominated & 0.0019 & 0.0214 & 0.0018 \\
Mixed regime (single run) & Mixed & 0.0291 & 0.0218 & 0.1327 \\
Rigor (8-seed mean) & Mixed & 0.0320 & 0.0210 & 0.1338 \\
Regression (covariate shift) & Harmful-dominated & ${\approx}0$ & 0.0864 & ${\approx}0$ \\
\bottomrule
\end{tabular}
\end{table}

\begin{figure}[ht]
\centering
\includegraphics[width=0.62\textwidth]{fig_regret_summary.png}
\caption{Regret-to-oracle across all sub-experiments of this chapter.  \KGA\ (blue)
  dominates always-adapt (orange) in the harmful and mixed regimes while remaining
  within range in the helpful regime.  Always-freeze (grey) is safe only in
  harmful-dominated regimes.  Source: aggregated from
  \texttt{knowability\_results.json}, \texttt{kbound\_harmful\_results.json},
  \texttt{mixed\_regime\_results.json}, \texttt{rigor\_multiseed.json},
  \texttt{regression\_covariate.json}.}
\label{fig:exp1-regret-summary}
\end{figure}

\subsection{Interpretation}

The synthesis table reveals the expected pattern from the theory of
Chapters~\ref{ch:theory1}--\ref{ch:theory2}: no fixed policy is universally best, and
\KGA\ adapts gracefully across regimes.

\begin{itemize}
  \item \textbf{Helpful-dominated (clean suite)}: always-adapt is near-oracle; \KGA\
    carries modestly more regret ($0.0167$ vs $0.0057$) because it abstains on tasks
    where adaptation would have helped.
  \item \textbf{Harmful-dominated (harmful fusion, regression)}: \KGA\ matches
    always-freeze and the oracle; always-adapt pays large regret ($0.0214$, $0.086$).
  \item \textbf{Mixed}: \KGA\ records lower regret than always-freeze ($0.029$ vs
    $0.133$) in-fold while remaining competitive with always-adapt ($0.022$); we read this
    row as a diagnostic, not a beats-both, since the mixed-stream comparison was withdrawn
    pending an out-of-fold re-run (Section~\ref{sec:exp1-mixed}).
\end{itemize}

The certificate's clearest evidence in this chapter is the suppression of harmful
adaptation: the $11\times$ regret reduction over always-adapt in the harmful-fusion
experiment (Section~\ref{sec:exp1-harmful}) and the matched oracle performance in the
regression experiment (Section~\ref{sec:exp1-regression}), in both of which \KGA\
\emph{matches} the safe baseline rather than beating both policies.  The corruption-grid
beats-both result --- where \KGA\ exceeds both trivial policies --- is the headline of
Chapter~\ref{ch:exp2}, to which this breadth chapter hands off.

% ============================================================
\section{Threats to Validity}
\label{sec:exp1-threats}
% ============================================================

\paragraph{Testbed specificity.}
All 123 tasks use precomputed anomaly scores from a fixed set of six detectors.  The
certificate's evidence vector $Z$ is tuned to this domain: it includes
\texttt{disagree}, the inter-detector disagreement feature that Theorem~\ref{thm:disagree}
and the ablation (Section~\ref{sec:exp1-ablations}) identify as the most informative
signal.  Results may not transfer directly to domains where no ensemble disagreement is
available.

\paragraph{Controlled corruption vs real sensor failure.}
The detectable-failure and covert-failure conditions in the mixed regime
(Section~\ref{sec:exp1-mixed}) are constructed via deterministic permutation and
distribution injection.  Real sensor or distribution failures may not match these
idealized forms; in particular, real covert failures may produce more or less
non-identifiability than the permutation construction.

\paragraph{Mild benefit in the clean condition.}
The clean condition has mean $B = -0.041$ (field
\texttt{true\_mean\_benefit\_B\_by\_condition.clean}): adaptation is mildly harmful, not
catastrophically so.  This makes always-adapt near-competitive in the mixed regime, and is
one reason the mixed-stream AUROC comparison is read here as a diagnostic rather than a
win.  The harmful-fusion experiment (Section~\ref{sec:exp1-harmful}) covers the severely
harmful case, but the mixed regime does not.  The decisive test --- catastrophically
harmful adaptation in a real neural network, and the corruption-grid setting where the
certificate exceeds both trivial policies --- is the headline application of the next
chapter, Chapter~\ref{ch:exp2}, to which this breadth chapter now hands off.

\paragraph{Certificate width vs coverage.}
At the canonical $\alpha = 0.10$, the conformal band is wide enough that only $35\%$ of
tasks receive a commit decision in the mixed regime.  In settings with small test batches
(batch size $16$ reduces coverage to $2\%$; Table~\ref{tab:exp1-batch}), the certificate
is extremely conservative and effectively defers nearly all decisions.  This is
information-theoretically honest but operationally limiting.

\paragraph{Regression experiment scale.}
The regression sub-experiment (Section~\ref{sec:exp1-regression}) uses 72 synthetic
instances with a known DGP.  It demonstrates cross-domain transfer but does not constitute
a large-scale validation.  Real covariate-shift regression tasks with estimated rather than
known density ratios may yield different results.

\paragraph{Discrepancies with kbound.tex.}

We flag the following numerical differences between the experiment protocol text in
\texttt{kbound.tex} (Section~7) and the JSON source files:

\begin{enumerate}
  \item \textbf{Harmful regime regret ratio.}  The paper states ``regret-to-oracle falls
    from $0.0214$ to $0.0019$ ($\approx 11\times$)''; the JSON files give
    always-adapt regret $= 0.0214$ and \KGA\ regret $= 0.0019$
    (\texttt{kbound\_harmful\_results.json}), confirming the numerics.
    However, the paper's framing compares \KGA\ ($0.7479$) to always-freeze ($0.7481$)
    and calls them equal; the JSON shows \KGA's AUROC is $0.7479$ and
    always-freeze is $0.7481$---a $0.0002$ gap that the paper rounds to zero.  This
    is not a contradiction but a rounding; we report the unrounded values in
    Table~\ref{tab:exp1-harmful}.

  \item \textbf{Mixed-regime AUROC: paper says KGA = 0.690, JSON says 0.6899.}
    The paper quotes $\text{AUROC} = 0.690$ for \KGA\ in the mixed regime; the JSON
    field \texttt{mean\_auc\_policies.K\_Bound} $= 0.6899$.  This is a rounding
    difference; we use the JSON value.

  \item \textbf{Rigor AUROC: paper says KGA $0.686 \pm 0.003$, JSON gives $0.6860 \pm 0.0033$.}
    Minor rounding; JSON values used here.

  \item \textbf{Rigor $t$-test: paper says $t = 62.2$, $p = 7.3 \times 10^{-11}$;
    JSON gives $t = 62.236$, $p = 7.262 \times 10^{-11}$.}  Minor rounding; JSON
    values used.

  \item \textbf{Regression decision counts.}  The paper says ``0 adapt, 11 freeze,
    61 abstain'' (total $= 72$) and the JSON confirms these exactly
    (\texttt{decision\_counts}).  No discrepancy.

  \item \textbf{Abstain $|\Delta|$ values.}  The paper says ``mean $|\Delta| = 0.021$
    on abstained vs $0.132$ on acted tasks''; the JSON gives $0.02150$ and $0.13237$
    respectively (\texttt{knowability\_results.json}).  Minor rounding; JSON values
    used.
\end{enumerate}

None of these discrepancies affect any qualitative conclusion; all result from rounding in
the paper's narrative.  The JSON files are the authoritative source for this monograph.
